\section{Rocq Mechanization}
\label{sec:rocq-mechanization}

In this section, we assess our Rocq formalization of mailbox types and Pat.
We start with a general overview of the mechanization's structure, and then move into the specifics of different aspects.

Figure \ref{fig:rocq-modules} shows the Rocq project's structure for mechanizing Pat.
The definition of types is required by the syntax module since it also defines what a Pat program is.
Substitution is a syntactic operation, however, the module also includes proofs about substitution on well-typed terms, thus depending on the typing rules.

The table in Figure \ref{fig:rocq-modules} shows the lines of code, number of definitions and lemmas for each file.
Module \coqe{Environments} contains, by far, the most amounts of lemmas.
This reflects the complexity stemming from environment combinations and environment subtyping.

The mechanizations of messages, mailbox patterns and mailbox types are straightforward.
Mailbox configurations mostly rely on the definitions of permutations of lists, already provided by Rocq's standard library.
% Thus, a significant amount of proofs simply make use of already existing lemmas, with only little adjustments needed.
The formalization of mailbox patterns requires more effort.
Defining the semantics of mailbox patterns as a relation between configuration and pattern, instead of possibly infinite multisets,
yield a mechanization that is adequate to work with \todo{ES}{Good?}.
The proofs about pattern equivalence and inclusion mostly follow the same approach of performing induction over a mailbox pattern.
While a substantial amount of lemmas need to be proven before being able to proof balancing (Lemma \ref{lemma:balancing}), most of
the proofs are manageable in size and complexity.
Since the difference in mailbox types between Pat and the mailbox calculus are mostly the usage annotations provided by quasi-linear types,
a mechanization of the mailbox calculus could certainly reuse our definitions and properties.

Formalizing environments is, by far, the most complex part of this mechanization.
% Even finding an acceptable representation took time and testing.
The library \emph{dblib} reduces some of the complexity around environments by proving operations and properties on them.
However, reasoning about environment combinations and environment subtyping still is far from simple.
Independent of the representation of environments that were considered, environment splitting and subtyping require a substantial amount of technical lemmas.
While most of the proofs of these properties follow the same general approach, they are lengthy and tedious to perform.
It may be possible that other representations of environments could reduce the complexity of such proofs.

% Before using \emph{dblib}, one approach was the usage of list concatination to reason about inserting values into an environment.
% For example, instead of writing $\envInsert{x}{A}{\Gamma}$ for inserting a type into $\Gamma$,
% one would have to specify $\Gamma = \Gamma_1 \, {++} \, (A, \Gamma_2$), where the length of $\Gamma_1$ is equal to $x$.
% Of course, specifying insertion in this way is tedious and error prone, as well as hard to reason about.
% The library \emph{dblib} provides a much easier and more understandable way, and reduces some of the complexity around
% environments by proving operations and properties on them.
% However, reasoning about environment combinations and environment subtyping still is far from simple.
% Independent of the representation of environments that were considered, environment splitting and subtyping require a substantial amount of technical lemmas.
% While most of the proofs of these properties follow the same general approach, they are lengthy and tedious to perform.
% It may be possible that other representations of environments could reduce the complexity of such proofs.

The mechanization of the syntax and typing rules of Pat is close to their pen-and-paper counterparts,
with the exception of using de Bruijn indices.
Defining the syntax and rules poses no problem after settling on the representation of variables and environments.
While the choice of using de Bruijn indices is reasonable, it comes with the downside of extensive reasoning about indices and lifting.
This is especially evident in the proofs required for the substitution lemma (Lemma \ref{lemma:substitution}).
While \emph{dblib} alleviates some of this burden, a substantial amount of manual work is still required.
Exploring whether other variable representations require less work is definitely worth to consider.

% \begin{table}[t]
%   \begin{center}
%     \begin{tabular}{|l | c | c | c|}
%       \hline
%       Module & LOC & Definitions & Lemmas\\
%       \hline\hline
%       Messages & 177 & 7 & 9\\
%       Mailbox Patterns & 828 & 8 & 44\\
%       Types & 685 & 35 & 29\\
%       Environments & 2658 & 18 & 110\\
%       Util & 48 & 6 & 2\\
%       Syntax & 180 & 16 & 0\\
%       Typing Rules & 302 & 4 & 11\\
%       Substitution & 941 & 2 & 26\\
%       Runtime & 959 & 18 & 35\\
%       \hline
%               & 6778 & 114 & 266\\
%       \hline
%     \end{tabular}
%   \end{center}
%   \caption{Table showing the lines of code, amount of definitions and lemmas in each file of the Rocq project.}
%   \label{table:rocq-modules}
% \end{table}

Due to \emph{dblib}, the substitution operation is provided for us with little effort.
% As a user, one simply has to provide a \emph{traverse} function over terms, counting how many binders are being passed.
That said, the lemmas required for the substitution lemma are quite technical and are mostly about a relation between two variables.
Since variables are considered values in Pat, reasoning about substituting one variable for another is required.
While \emph{dblib} provides tactics for automatically proving some of these properties, it often would not be able to do so,
requiring manual proofs.
Additionally, during the process of proving the substitution lemma, it was discovered that the proof is not as easy as it first seems.
Due to linear typing and environment combinations, one has to consider more cases than one would expect.
Thus, the substitution lemma's proof is by far the most complex and largest proof in this mechanization.

% Further, we successfully initiated the mechanization of Pat's operational semantics.
% While it seems straightforward, a big hurdle are in fact de Bruijn indices, representing variables.
% One complex aspect in that regard is the definition of structural congruence.
% It allows for the rearranging of name restrictions, and in turn, requires operations that change de Bruijn indices.
% For example, for some term $(\nu \, x)(\nu \, y)P$, with structural congruence, we are able to swap both restrictions, resulting in
% $(\nu \, y )(\nu \, x)P$.
% While this does not require any modifications when using classical names as variables,
% for de Bruijn indices, we have to exchange the indices in $P$.
% Additionally, whenever a new mailbox name is created dynamically, it is assigned index $0$.
% Thus, we have to lift all indices in all current configurations by one.
% This constant manipulation of indices substantially increases the complexity of proving properties and
% even formulating appropriate definitions, such as structural congruence.
% While it is not yet clear how much work is required to formalize the operational semantics of Pat,
% it certainly is not an easy task.

During the mechanization, multiple oversights and inaccuracies in the original definitions and proofs were discovered.
For instance, the rules for environment subtyping and pattern residuals, as presented in the paper \cite{fowlerSpecialDeliveryProgramming2023}, are missing cases.
Additionally, while trying to prove some of the properties claimed, the proofs seemed more difficult than initially expected.
The proof of the substitution lemma (Lemma \ref{lemma:substitution}), for example, did not account for the complexity of environment combinations in the
case of let-bindings.
This resulted in examining the proof at a greater detail, requiring additional intermediate lemmas.
A similar problem was discovered while trying to prove the preservation property.
In some cases, subtyping was not accounted for, skipping over some of the most complex parts of the proof, as it turned out.
With the help of a rigorous formalization in Rocq, we were able to identify these oversights.
All found issues have been acknowledged and consequently fixed by the authors of Pat in an updated version of their original paper \cite{fowlerSpecialDeliveryProgramming2025}.
